\section{A finite support-cover criterion}

\subsection*{Finding F-018: two monomial order ideals and one linear edge}
Let \(\Gamma\subseteq\mathbb Z_{\geq0}^{r-1}\) be a
nonempty finite coordinatewise downward-closed set. Write
\(E=\{0,e_2,\ldots,e_r\}\). Call \((P,Q)\) admissible
if \(P,Q\subseteq\Gamma\) are nonempty downward-closed
sets and \(\Gamma\subseteq P+Q+E\). Define
\begin{equation}\label{eq:F018kappa}
 \kappa(\Gamma)=\min_{(P,Q)\ \mathrm{admissible}}
                                      \bigl(|P|+|Q|-1\bigr).
\end{equation}
Here the sum is the set of sums of multi-indices. The minimum
exists, since \(P=\Gamma,Q=\{0\}\) is admissible, and
\[
                 \max_{\gamma\in\Gamma}|\gamma|
                         \leq\kappa(\Gamma)\leq|\Gamma|.
\]
This is a finite optimization problem for this particular
graph construction. No polynomial-time algorithm or optimality
among all linear determinant representations is asserted.

\paragraph{Universal identity.}
For any integer \(h\geq\max_{\gamma\in\Gamma}|\gamma|\),
put
\[
               f(L,z)=\sum_{\gamma\in\Gamma}
                         c_\gamma z^\gamma L^{h-|\gamma|}.
\]
Every admissible pair \(P,Q\) yields a homogeneous linear
matrix with \(S=|P|+|Q|-1\) vertices and
\begin{equation}\label{eq:F018det}
                 \det M=\sum_{\gamma\in\Gamma}
                         c_\gamma z^\gamma L^{S-|\gamma|}.
\end{equation}
In particular, if \(n\geq\max\{h,\kappa(\Gamma)\}\),
every form \(L^{n-h}f\) in this coefficient space is an
\(n\)-by-\(n\) homogeneous linear determinant.

\begin{proof}
Choose one decomposition
\(\gamma=\alpha_\gamma+\beta_\gamma+\epsilon_\gamma\)
with \(\alpha_\gamma\in P\), \(\beta_\gamma\in Q\), and
\(\epsilon_\gamma\in E\), for every \(\gamma\).
Make the two monomial trees on \(P\) and \(Q\), using
the parent rule which deletes the last nonzero coordinate.
Downward closure guarantees all parents and root paths.
Orient the first tree away from its root and the second
toward its root, merge the two roots, and keep all other
vertices distinct. Label tree edges by the negative of the
corresponding variable.

Insert one coefficient-edge summand from \(\alpha_\gamma\)
to \(\beta_\gamma\), labeled \(c_\gamma L\) if
\(\epsilon_\gamma=0\), or \(c_\gamma z_i\) if
\(\epsilon_\gamma=e_i\). Add parallel summands in their
matrix entry. Put an \(L\)-loop at every noncentral
vertex; the central diagonal has only its coefficient labels.
Every non-loop cycle contains the central vertex and exactly
one coefficient edge. Its two paths are unique.

For \(\epsilon_\gamma=0\), the cycle has
\(|\gamma|+1\) vertices; for a unit vector it has
\(|\gamma|\) vertices. The determinant sign cancels the
negative tree-edge signs in both cases. The chosen coefficient
edge and the unused loops give precisely
\(c_\gamma z^\gamma L^{S-|\gamma|}\). This includes
central coefficient loops. The vertices of a cycle are
distinct, so all exponents are nonnegative. This proves
\eqref{eq:F018det} and the bound \(S\geq|\gamma|\)
for each \(\gamma\). One can also see the latter directly
from \(|P|\geq|\alpha_\gamma|+1\) and
\(|Q|\geq|\beta_\gamma|+1\).
Padding by \(L I_{n-S}\) gives the assertion at size \(n\).
All identities are polynomial identities over the integers
in the coefficient parameters.
\end{proof}

\paragraph{Two actual multiplicity criteria.}
For a partition \(\lambda\) of length \(r\geq2\), define
\[
 \Gamma_{h,\lambda}=
 \left\{\gamma\in\mathbb Z_{\geq0}^{r-1}:|\gamma|\leq h,
      \quad\sum_{i=j}^r\gamma_i\leq\lambda_j
                                      \ (2\leq j\leq r)\right\}.
\]
For internal padding, \(n>m\geq2\), and
\(\lambda\vdash nd\), one has
\begin{equation}\label{eq:F018q}
 n\geq\kappa(\Gamma_{m,\lambda})
       \quad\Longrightarrow\quad q^m_\lambda(d[n])\leq K_n(\lambda).
\end{equation}
If \(r\leq\min\{d,n^2\}\) and \(b=\lambda_2\), then
\begin{equation}\label{eq:F018K}
 n\geq\kappa(\Gamma_{b,\lambda})
       \quad\Longrightarrow\quad K_n(\lambda)=a_\lambda(d[n]).
\end{equation}
The first comparison uses the permanent degree; the second
uses the stabilized inner degree. They must not be interchanged.

\begin{proof}
For \eqref{eq:F018q}, discard the immediate case \(q=0\).
Positive \(q\) implies \(r\leq\min\{d,m^2\}\).
Take any determinant-ideal highest-weight vector of weight
\(\lambda\). Align the padding linear form with the first
coordinate on the dense open set used in M-009. The support
theorem F-016 permits deleting exactly the coefficient monomials
outside \(\Gamma_{m,\lambda}\) without changing its value.
The universal identity above makes the retained padded form
a determinant at size \(n\). Consequently that vector
vanishes on the permanent orbit and its reduced closure.
This holds for the whole highest-weight space, so it gives
the ideal multiplicity-space inclusion and the inequality
in \eqref{eq:F018q}, as in F-017.

For \eqref{eq:F018K}, the set \(\Gamma_{b,\lambda}\)
contains the pure second-variable power of degree \(b\).
Hence \(\kappa(\Gamma_{b,\lambda})\geq b\), and
the hypothesis implies \(n\geq b\).
F-014 stabilizes the ambient multiplicity at
\(\mu=(db-|\bar\lambda|,\bar\lambda)\vdash db\).
F-016 puts the entire highest-weight space of weight \(\mu\)
in the coefficient algebra of \(\Gamma_{b,\lambda}\).
Its dual padded coefficient space is contained in the
determinant closure by \eqref{eq:F018det}. The diagonal,
invertible padding pullback of F-015 preserves the independence
of all copies. Thus \(K_n(\lambda)\geq a_\lambda(d[n])\);
the reverse ambient quotient inequality proves equality.
\end{proof}

The preceding exact two-tree constructions give special choices of
the two order ideals and the coefficient-edge decomposition
after intersecting their vertex sets with \(\Gamma\), if
necessary. Every used endpoint and its parents are coordinatewise
bounded by their supported monomial, so this pruning retains
the required decompositions and can only reduce the size.
In applying \eqref{eq:F018q}, a pair may encode the entire
nested profile, instead of replacing it by just one row cut.
The elementary bound \(\kappa(\Gamma)\leq|\Gamma|\)
already becomes useful when joint terminal constraints make
the full support much smaller than every one-cut enlargement.

\paragraph{The permanent version preserves its distinguished padding entry.}
For \(r\leq\min\{d,m^2\}\), \(b=\lambda_2\), and
\(n>m\geq2\), there is also the exact statement
\begin{equation}\label{eq:F018per}
 m\geq\kappa(\Gamma_{b,\lambda})
       \quad\Longrightarrow\quad
          q^m_\lambda(d[n])=K_n(\lambda)=a_\lambda(d[n]).
\end{equation}
This determines full permanent-side multiplicities in additional
families. It is not a new determinant-side exclusion when
\eqref{eq:F018K} already applies.

\begin{proof}
In the graph for an admissible pair, replace each negative
tree-variable label by its positive counterpart. A permanent
cycle cover has no permutation sign, so its contribution is
now exactly the same monomial as in \eqref{eq:F018det}.
Thus the resulting matrix \(A\) satisfies
\[
                 \operatorname{perm}A
                    =\sum_{\gamma\in\Gamma}
                              c_\gamma z^\gamma L^{S-|\gamma|}.
\]
If \(S\geq2\), every noncentral diagonal entry is exactly
\(L\); coefficient edges never create a noncentral diagonal
because the two tree copies are distinct. A simultaneous
row-and-column permutation moves one such entry to position
\((1,1)\), without changing the permanent. Adjoin
\(L I_{m-S}\) to reach size \(m\).
If \(S=1\), first adjoin one \(L\)-diagonal entry and
move it to position \((1,1)\); this is possible because
\(m\geq2\). The remaining padding is the same.

Apply this to \(\Gamma=\Gamma_{b,\lambda}\), so that
\(b\leq S\leq m\). For every form in its degree-\(b\)
coefficient space the constructed matrix has
\[
                 A_{11}=L,\qquad
                 \operatorname{perm}A=L^{m-b}f.
\]
Substitution in the internally padded permanent gives
\(A_{11}^{n-m}\operatorname{perm}A=L^{n-b}f\).
Every such substitution lies in its orbit closure by density
of invertible endomorphisms. The same all-copy stabilization
and invertible padding restriction used above now give
\(q^m_\lambda\geq a_\lambda\). The ambient quotient
bound forces equality. Since \(n>m\geq\kappa\),
\eqref{eq:F018K} also gives \(K_n=a_\lambda\).
\end{proof}

\paragraph{Scheme and lambda-operation interpretation.}
The order ideals are finite monomial supports in a torus-graded
symmetric-power module, under SGA~3 I4.7.3. Their coefficient
schemes map to schemes of linear matrices over \(\mathbb Z\),
with determinant and permanent identities just proved.
The highest-weight support is supplied by exterior minors
and Cauchy/Schur functors (Macdonald I(4.3), Appendix~A(5.3)--(5.6);
SGA~6 VI3.3). The coordinate-ring conclusions retain the
specific quotient and evaluation maps: this is exactly the
additional data missing from a bare effective lambda class
in F-010. Characteristic-zero factorial invertibility and
semisimplicity are used only in the all-copy lift. The graph
identities themselves are integral. No saturation or general
positive effective Burnside lift is deduced.
